% Build the layer, Chapter 6, first build: one escrow for both sides of a job.
% Inserted at the end of \S\ref{sec:float-plan}, after the table of
% settlement outcomes. It builds the daemon's transaction of
% Protocol~\ref{prot:escrow} and the settlement of the section's worked
% example, as arithmetic in the kernel's decision core.
%
% No code or output in this file is typed. \pdexcerpt, \pdrun and
% \pdrunlines set files under excerpts/, which the kernel's
% tools/excerpts.mjs cuts from its source and its recorded runs (see
% figures/pd-listing.tex). The figure draws one of those recorded tables as
% it was printed. To change a listing or the figure, change the kernel and
% cut again. The prose names functions, never line numbers, so that a
% listing which moves in its file does not make a sentence wrong.
%
% \pdlistingroom{n} before a listing of n lines keeps it on one page. When a
% listing changes length, change the number with it.
\subsection{Build the layer: one escrow for both sides of a job}
\label{sec:he-plan}

\noindent\emph{Express lane: the middle of the escrow ceremony is built
twice. The first keeper checks each amount and works out each share by
itself. It gets this section's worked settlement right, and then it
overdraws a wallet and loses a unit. The second reserves both amounts in
one decision, and takes the second part of every split by subtraction. A
walk takes it through 4,213 states, and it pays one plan out 600 ways. A
solver does the rest with one line of arithmetic taken out of its hands,
and the section says which line and why.}

Protocol~\ref{prot:escrow} has a transaction in the middle. The daemon
checks the balances again and debits the bounty and the bond into escrow,
in one commit. Later a settlement splits each of the two by the share of
the work that was done. \pdkernelname{} builds those two moves here. The
signatures on either side of them wait for a later build, since the cards
of \pdchapref{anchor}{The Anchor Protocol} are the means to make them.

Money is counted in whole units, so every amount is a whole number. A
\emph{keeper} holds the books: a wallet for each principal and the
commons, and every plan with the two amounts it has in escrow.

\subsubsection{Naive: each amount by itself}

The obvious keeper follows the ceremony's words. To open a plan it checks
that the requester's wallet covers the bounty and that the provider's
covers the bond, and debits both. To pay a plan out it follows
Table~\ref{tab:settlement}: each side's share is its proportion of the
amount, rounded down to whole units (Listing~\ref{alg:he-plan-apiece}).

\pdlistingroom{25}
\pdexcerpt[caption={The first try. Each wallet is checked for its own
amount, and each share is worked out by
itself.},label={alg:he-plan-apiece}]{ch06-apiece}

This keeper is not careless. It refuses a plan that a wallet cannot cover,
and it pays a plan out once. On this section's own example it is right
(Listing~\ref{alg:he-plan-run-partial}).

\pdlistingroom{15}
\pdrun[caption={The reduced contract of \S\ref{sec:float-plan}, run on the
first try. Nobody here is asked who they are: \texttt{both} and
\texttt{judge} say who would
ask.},label={alg:he-plan-run-partial}]{partial-apiece}

Bob ends with 240 of his 400 units and Alice with 200, and the other 60 are
in the commons. Those are the net changes in the example's table. A reader
who checked this keeper against the chapter would pass it.

\subsubsection{Break: one wallet behind two amounts}

A principal can be both sides of a plan. Bob asks for a piece of work and
Bob's own agent does it, so the bounty and the bond come out of one wallet.
The ceremony has a clause for this case: a wallet used for both must cover
the two amounts together (Listing~\ref{alg:he-plan-run-wallet}).

\pdlistingroom{11}
\pdrun[caption={The break. A wallet of 400 stands behind a bounty of 400
and a bond of 100.},label={alg:he-plan-run-wallet}]{one-wallet-apiece}

Step 2 passed both checks. The wallet holds 400, which covers a bounty of
400. It holds 400, which covers a bond of 100. Neither check knows of the
other, and the wallet ends 100 below nothing. Escrow says 500 where 400
units exist. A bond is meant to be something the provider stands to lose
beyond the bounty. This one is the bounty's own units, counted a second
time.

\subsubsection{Break again: a unit in nobody's hands}

The second fault needs two principals and amounts that the milestones do
not divide (Listing~\ref{alg:he-plan-run-thirds}).

\pdlistingroom{15}
\pdrun[caption={The second break. A plan of 110 units is paid out, and 108
arrive.},label={alg:he-plan-run-thirds}]{thirds-apiece}

One milestone of 3 was done. A third of the bounty is 33 units and a third
of a unit, so Alice earns 33. Two thirds is 66 and two thirds of a unit, so
Bob has 66 back. Each share is right to within a unit, and together they
come to 99. The bond goes the same way, with 3 back and 6 forfeited out of
10. Two of the 110 units are now nowhere in the books. Design
Invariant~\ref{prop:conservation} is false after step 4, and no request
was refused on the way.

The example above did not show this because its numbers are round. A
bounty of 400 at 2 milestones of 5 divides evenly, and so does a bond of
100. A worked example with round numbers tests the rule it illustrates and
little of the arithmetic under it. The verification below pays one plan
out 600 ways, and Figure~\ref{fig:he-odd-unit} shows how much the round
numbers hid.

\subsubsection{Fix: decide the whole, and subtract}

The two faults are one: the keeper checked parts where the rule was about
a whole. \pdkernelname{} puts both moves in the decision core
(\weld{prereq:sec:prereq-build}{decision core}), which is shown both of a
plan's wallets at once (Listing~\ref{alg:he-plan-reserve}).

\pdlistingroom{41}
\pdexcerpt[caption={The decision to open a plan. \texttt{one\_wallet} is
its guard.},label={alg:he-plan-reserve}]{ch06-reserve}

\texttt{covered} is a subtraction that refuses when it would go below
nothing, and the question mark after it hands that refusal back before
anything is built. For two wallets each is asked for its own amount. For
one wallet, \texttt{one\_wallet} sends the decision down the other arm,
where the wallet is asked for the sum. Either way the answer is both
wallets as they will be, with a plan that holds both amounts, or it is a
refusal. The kernel stores the answer in one step. Bob's wallet of 400 is
asked for 500 (Listing~\ref{alg:he-plan-run-wallet-core}).

\pdlistingroom{11}
\pdrun[caption={The fix. One wallet is asked for 500 and holds
400.},label={alg:he-plan-run-wallet-core}]{one-wallet-core}

A split subtracts (Listing~\ref{alg:he-plan-split}).

\pdlistingroom{21}
\pdexcerpt[caption={A split. The second part is what the first
leaves.},label={alg:he-plan-split}]{ch06-split}

The two parts are the amount whatever \texttt{share} returns, because the
second is defined as the rest. Rounding can no longer lose a unit. It can
only decide who has it. Here the provider's share is rounded down, so the
odd unit of a bounty goes back to the requester and the odd unit of a bond
goes to the commons. The milestone rule of this chapter does not say where
an odd unit goes, and this is \pdkernelname{}'s choice: a provider is
never paid above its share of the work. Exercise~\ref{ex:he-plan-odd-unit}
weighs the other choice.

The subtraction must not go below nothing, and that rests on
\texttt{counted}, which refuses a count of more milestones done than there
are. When \texttt{done} is no more than \texttt{of}, the amount times
\texttt{done} is no more than the amount times \texttt{of}, and dividing
both by \texttt{of} gives a share that is no more than the amount. That
argument comes back below. A settlement is then two splits by one count
(Listing~\ref{alg:he-plan-settle}).

\pdlistingroom{30}
\pdexcerpt[caption={The decision to pay a plan out. Each amount is split
by the same count, and each part has one place to
go.},label={alg:he-plan-settle}]{ch06-settle}

\texttt{open} refuses a plan that was paid out earlier, so the escrow
leaves once. The plan of Listing~\ref{alg:he-plan-run-thirds} now pays out
all 110 of its units (Listing~\ref{alg:he-plan-run-thirds-core}).

\pdlistingroom{6}
\pdrunlines[caption={The fix. The odd unit of the bounty is Bob's, and the
odd unit of the bond is the
commons'.},label={alg:he-plan-run-thirds-core}]{thirds-core}{1,11-15}{1 and 11\textendash 15}

\keyidea[Check the whole, and derive the last part]{Each amount was
covered and each share was right to within a unit, and both times the
whole was wrong. A rule about a sum has to be checked on the sum. Where
parts must come to a whole, work out all of them but one and subtract, so
that no way of rounding can leave the books short.}

\Needspace{12\baselineskip}
\subsubsection{Verify: small cases in full, and then any numbers}

\paragraph{Every state.}
A walk tries everything a keeper can be asked, from every state it can
come to, within small bounds: 2 principals, deposits of 1 or 2 units up to
3 in all, a bounty and a bond of up to 2 units, counts of up to 3
milestones, and 2 plans. Most of what it tries should be refused. A plan is asked for with no bond, or from a
wallet that holds too little. A plan is paid out twice, or by a count of
4 milestones done out of 3. Each answer is held to a second set of rules,
which are the ceremony's clause and the settlement table written out on
plain numbers in a file that shares no code with the core. The bounds are
small because states multiply, and the walk has to run each time the
kernel's tests do.

The kernel's keeper reaches 4,213 states. It makes all 20,120 moves that
the rules allow and none of the 164,482 they refuse. The first try takes
650 moves that the rules refuse and comes to 682 states with a wallet
below nothing. Of its payouts 6,918 are short, and they lose 8,560 units
between them (Listing~\ref{alg:he-plan-states}).

\pdlistingroom{18}
\pdrun[caption={Every state each keeper can reach, held to this section's
own rules.},label={alg:he-plan-states}]{plans}

\paragraph{Every count.}
The walk is where one wallet meets two amounts and where two plans meet.
With amounts of 2 units it sees little rounding. So \pdkernelname{} also
pays one plan out at every count of 2 to 6 milestones, short of all of
them, for every bounty from 1 to 40 units. That is 600 payouts by each
keeper (Figure~\ref{fig:he-odd-unit} and
Listing~\ref{alg:he-plan-splits}).

% Brief (grammar: a grid of every case tried, with a mark on each case that
%   went wrong; atlas IV/fig:he-odd-unit).
% Reader question: when does a payout lose a unit, if each side's share is
% computed separately and rounded down?
% Claim: whenever the milestones do not divide the bounty times the count
% done. That is 424 of the 600 payouts tried, which the picture counts as
% it draws them, and no bounty up to 40 units is exact at every count.
% Evidence: one plan paid out by the first ledger at every count of 2 to 6
% milestones, with at least one not done, for every bounty from 1 to 40
% units [internal, harbor splits --each; session splits-each.txt, drawn as
% printed, and tests/harbor-research/test_he_plan_numbers.py compares the
% counts in the picture and the caption with it].
% Must distinguish: a payout that lost a unit from one that lost none;
% one count of milestones from the next.
% Grammar chosen: the bounty along the bottom and the count down the side,
% so that the rule shows as stripes whose width differs by row, and a
% reader can find any one payout.
% Grammar rejected: a bar for each number of milestones, which gives the
% proportion that lost a unit and hides the rule; the table itself, which
% is 600 rows.
% Five-second test: is any column empty?
\begin{figure}[H]
\centering
\pdfigurehue{pdcobalt}%
\begin{tikzpicture}[pd figure,x=.232cm,y=.3cm,
    lost/.style={fill=pderror},
    exact/.style={fill=pdinkmuted!70}]
  % The table is the recorded run, read as it was printed: the first line
  % is the command, and every remark starts with '#'.
  \pgfplotstableread[skip first n=1]{excerpts/splits-each.txt}\pdoddunit
  % Counts taken from the table as it is drawn.
  \ifdefined\pdoddlost\else \newcount\pdoddlost \newcount\pdoddall \fi
  \global\pdoddlost=0 \global\pdoddall=0
  % A row for each count, with a gap between one number of milestones and
  % the next: the row of `done` of `of` milestones.
  \newcommand{\pdoddrow}[2]{-((#2-2)*(#2-1)/2+#1-1)-.6*(#2-2)}
  \pgfplotstablegetrowsof{\pdoddunit}
  \pgfmathtruncatemacro{\pdlastsplit}{\pgfplotsretval-1}
  \foreach \i in {0,...,\pdlastsplit}{
    \pgfplotstablegetelem{\i}{of}\of\pdoddunit \let\mi\pgfplotsretval
    \pgfplotstablegetelem{\i}{done}\of\pdoddunit \let\dn\pgfplotsretval
    \pgfplotstablegetelem{\i}{bounty}\of\pdoddunit \let\bt\pgfplotsretval
    \pgfplotstablegetelem{\i}{lost}\of\pdoddunit \let\lf\pgfplotsretval
    \pgfmathsetmacro{\y}{\pdoddrow{\dn}{\mi}}
    \global\advance\pdoddall 1
    \ifnum\lf>0 \global\advance\pdoddlost 1
      \fill[lost] (\bt-.41,\y-.36) rectangle (\bt+.41,\y+.36);
    \else \fill[exact] (\bt,\y) circle[radius=.6pt];\fi
  }
  \foreach \mi in {2,...,6}{
    \pgfmathtruncatemacro{\most}{\mi-1}
    \foreach \dn in {1,...,\most}{
      \pgfmathsetmacro{\y}{\pdoddrow{\dn}{\mi}}
      \node[pd label,text=pdinkmuted,anchor=east] at (.1,\y) {\dn\ of \mi};
    }
  }
  \pgfmathsetmacro{\pdoddfloor}{\pdoddrow{5}{6}-1.15}
  \foreach \bt in {1,5,10,15,20,25,30,35,40}{
    \node[pd label,text=pdinkmuted,anchor=north] at (\bt,\pdoddfloor) {\bt};
  }
  \node[pd label,anchor=north] at (20.5,\pdoddfloor-1.5) {the bounty, in units};
  \node[pd title,anchor=south west] at (.55,.75)
    {a unit is lost in \the\pdoddlost\ of \the\pdoddall\ payouts};
\end{tikzpicture}
\caption{\textbf{Two shares rounded down separately lose a unit unless
the milestones divide the bounty times the count done.} A red square is a
payout by the first ledger that lost one unit, and a dot is a payout that
lost none. [internal, \texttt{harbor splits -{}-each}]}
\label{fig:he-odd-unit}
\end{figure}


\pdlistingroom{16}
\pdrun[caption={One plan paid out 600 ways by each
keeper.},label={alg:he-plan-splits}]{splits}

A payout comes out whole only when the milestones divide the bounty times
the count done, which is why each row of the figure has stripes of its own
width. No bounty up to 40 is clear in every row. The kernel's keeper lost
no unit in any of the 600. Exercise~\ref{ex:he-plan-floors} asks for the
proof that two shares rounded by themselves lose one unit of an amount or
none.

\paragraph{Any numbers.}
The walk stops at amounts of 2 units and the sweep at 40. Kani lifts the
bound (\weld{anchor:sec:correct}{Kani}): it compiles the core's own
functions and asks a solver about every input at once. Every wallet and
every amount is 64 bits that the solver chooses, and so is each count.

One line stands in the way. \texttt{share} multiplies and divides, and
Kani's solver reasons one bit at a time, which is slow at both. With the
real \texttt{share}, one split took it three minutes for amounts of 16
bits, and a settlement with amounts of 32 bits was stopped after ten. So
the proof is in two parts. The harness on a settlement does not run
\texttt{share}. It runs \texttt{some\_share} in its place, a function
that returns any number at all that is no more than the amount, when the
count is one a plan can have. That is everything a settlement uses of a
share. A second harness then asks the real \texttt{share} whether it
keeps that promise, for numbers small enough for the solver
(Listing~\ref{alg:he-plan-kani-settle}).

\pdlistingroom{33}
\pdexcerpt[caption={A settlement as a Kani harness, and what the harness
is given in place of
\texttt{share}.},label={alg:he-plan-kani-settle}]{ch06-kani-settle}

The second line of the listing makes the swap. The harness then says what
each side can gain. The commons takes no more than the bond and the
requester has back no more than the bounty, and what those two do not get
is exactly the provider's gain. It speaks side by side because the solver
was stopped after ten minutes on one equation with every gain in it, and
takes 23 seconds this way.

Kani's verdicts are in Listing~\ref{alg:he-plan-kani}, and the fourth is
the other part. It runs the real \texttt{share}, for amounts and counts of
8 bits, and finds that a share is no more than the amount. So the tool has
checked two things. For any wallets and amounts of 64 bits, a settlement
pays out what was reserved, if a share is no more than the amount. For
amounts and counts up to 255, a share is no more than the amount. Past 255
that fact rests on the sentence of arithmetic in the fix above, and on no
tool. The last verdict is a control: it puts two shares rounded by
themselves to the same question, and fails as it must.

\pdlistingroom{11}
\pdrunlines[caption={Kani~0.68 on the core's rules for a
plan.},label={alg:he-plan-kani}]{kani-plans}{1-11}{1\textendash 11}

\paragraph{Each guard removed.}
This build has five guards, and the kernel has a build with each one taken
out. The walk is run on each (Listing~\ref{alg:he-plan-guards}).

\pdlistingroom{17}
\pdrun[caption={The walk again, on five builds of the core with one guard
removed from each.},label={alg:he-plan-guards}]{plans-guards}

Every build goes wrong in the walk, and every build fails some harness
under Kani (Listing~\ref{alg:he-plan-kani-guards}). The last row is Design
Invariant~\ref{prop:nospawn}, as far as a build that starts no worker can
take it: an open plan has a bond.

\pdlistingroom{15}
\pdrunlines[caption={The harnesses again, on the same five
builds.},label={alg:he-plan-kani-guards}]{kani-plans}{1,12-25}{1 and 12\textendash 25}

The first row of Listing~\ref{alg:he-plan-guards} repays a second look.
Without \texttt{one\_wallet} the core treats one wallet as two. It debits
each for its own amount and answers with two balances, and the kernel
keeps the one it is handed second. No wallet goes below nothing. The
bounty is in escrow and still in the wallet, and 1,596 states hold more
units than were deposited. That is more than the 636 wrong moves, because
a plan the wallet does cover is booked wrongly too. The first try showed
this fault as a debt, and this build shows it as units from nowhere.
\pdscope{The walk covers 2 principals and 2 plans with amounts of up to 2
units, within the bounds that Listing~\ref{alg:he-plan-states} prints, and
it holds each keeper to rules written by the same hand as the core. The
sweep covers one plan, bounties of 1 to 40 units and 2 to 6 milestones.
The harnesses cover one move from books that add up. For numbers past 255
they take it as given that a share is no more than the amount.}

\begin{pdboundary}
\setlength{\parskip}{4pt plus 1pt}
Nobody signs. The ceremony has both sides sign one plan, and has the
daemon sign a receipt for the worker. This build is the transaction
between those. Whoever reaches the keeper can open a plan between two
principals.

Nobody is asked who counts. A payout takes its count of milestones from
whoever asks, and a count of none forfeits the whole bond. The oracle of
Table~\ref{tab:settlement} is not built, and neither is the hold for a
dispute.

Every milestone weighs the same, and a plan has one requester and one
provider. The lease and the licence of \S\ref{sec:three-on-one} would be
more entries in the same escrow, and they are not here.

The books live in the keeper's memory, as Chapter~0's number does.
\pdchapref{bonded}{The Bonded Commons} builds a project's funds on the
same footing, and the two sets of books are not yet joined.
\end{pdboundary}
